% AI-produced by GPT-6 Astra (OpenAI, ChatGPT Pro mode). No priority claim.
\subsection{Descent through the generating cover at birth}

For a new object $P$ of the repaired C.1 construction, its
\emph{generating cover at birth} consists of its prescribed incoming
copy maps.  These are the two summand maps for a sum, or the maps
$\alpha_i:T\to P$ for a quasi-filter.  Every nonidentity arrow into
$P$ factors through a member of this cover.  For the original
object $W$, use the three singleton maps $v_i:V\to W$ as its
generating cover.  The original singleton $V$ will be treated
separately.

Two distinct points $p,q\in P$ are called a \emph{forbidden pair}
if no member of this cover contains both.  This terminology refers
only to the fixed source $P$ and its actual generating cover.

\begin{lemma}[Proved: routing a forbidden pair]
\label{C:birth-forbidden-pair}
Let $P\ne V$ be an object with its generating cover at birth, and
let $p,q\in P$ be a forbidden pair.  At every later stage, if
$h:P\to X$ and $g:Y\to X$ satisfy
\[
                         h(p),h(q)\in g(Y),
\]
then $h$ factors through $g$ by an arrow of the constructed
category.
\end{lemma}

\begin{proof}
We induct on the stage of $X$, starting at the birth stage of
$P$.  At that stage the only possible arrow with source $P$ is
its identity.  If the comparison arrow into $P$ is proper, it
factors through one generating image, which cannot contain the
forbidden pair.  The identity comparison is immediate.  There
are no arrows from $P$ to distinct objects born at the same stage,
or to earlier objects.  This also proves the base case $P=W$ in
the original category.

Let $X$ now be a new quasi-filter target, with old target $T$,
old pair $\phi_0,\phi_1:R\rightrightarrows T$, and generating maps
$\alpha_i:T\to X$.  The identity comparison is immediate.
Otherwise write
\[
                         h=\alpha_jh^*,\qquad
                         g=\alpha_ig^*.
\]
When $j\ne i$, choose the next index $j'$ on the shortest cyclic
arc from $j$ to $i$.  Its length is at most $2k\leq s-2n$.
Property (10) puts the two displayed points in the overlap of
copies $j,j'$.  Property (11) identifies that overlap with
\[
                       \alpha_j\phi_r(R)
                          =\alpha_{j'}\phi_t(R)
\]
for suitable $r,t\in\{0,1\}$.  Injectivity of $\alpha_j$ places
$h^*(p),h^*(q)$ in $\phi_r(R)$.  The old induction hypothesis,
applied to the same forbidden pair in the fixed source $P$,
supplies $l:P\to R$ with $h^*=\phi_rl$.  Thus the entire arrow
$h$ also factors through the next copy:
$h=\alpha_{j'}\phi_tl$.

Repeat along the chosen shortest arc.  Once $h=\alpha_i\widetilde h$,
injectivity places $\widetilde h(p),\widetilde h(q)$ in $g^*(Y)$.
The old assertion factors $\widetilde h$ through $g^*$, giving
the required factorization through $g$.

At a disjoint-sum target, any nonempty intersection puts the two
proper arrows in the same summand, where the old assertion
applies.  This completes the induction.  The argument uses the
actual prescribed incoming factorizations, not arbitrary graph
embeddings.
\end{proof}

\begin{lemma}[Proved: pairwise visibility has a common copy]
\label{C:birth-cover-helly}
Let $P\ne V$ have its generating cover at birth.  A nonempty
subset $U\subseteq P$ with no forbidden pair is contained in one
member of that generating cover.
\end{lemma}

\begin{proof}
For $W$, the assertion says that a set with at most one point is
contained in one singleton image.  For a disjoint sum, any pair
of points in different summands is forbidden, so all of $U$ lies
in one summand.

For a quasi-filter target, let
\[
                   I_x=\{i:x\in\alpha_i(T)\}
                              \subseteq\mathbb Z/s
\]
be the visibility set of a point.  The repaired coordinate
formulas, as recorded in Proposition~\ref{C:necklace-set-colimit},
make each $I_x$ either the full index cycle or a circular interval
with at most
\[
                              w=\max(1,2n-2)
\]
indices.  The extra $1$ covers private points even when $n=1$.
For every C.1 attachment,
\[
 \frac{s}{3n}\geq
 \left(1-\frac{a_{m+1}}{a_m}\right)^{-1}>4.
\]
The left side is an integer, so $s\geq15n>3w$.

The hypothesis on $U$ says exactly that its visibility sets meet
pairwise.  Full-circle sets impose no restriction.  If a proper
set occurs, select one and lift it to an integer interval
$I=[0,a]$ with $a\leq w-1$.  Every other proper visibility interval
has a unique lift meeting $I$.  Every such lift is contained in
$[-(w-1),a+(w-1)]$, whose span is at most $3w-3<s$.
Thus two of these lifts intersect whenever their circular images
do: distinct points in that common window cannot differ by $s$.
They are therefore pairwise intersecting ordinary intervals.
Their largest left endpoint is at most their smallest right
endpoint, so they have a common index.  That index also belongs
to every full-circle visibility set.  If all the sets were full
circles, any index would work.  Hence some generating copy
contains every point of $U$.
\end{proof}

\begin{proposition}[Proved: intersections are existing arrow images]
\label{C:arrow-image-intersections}
For any two arrows $h:P\to X$ and $g:Q\to X$ of a repaired C.1
stage, their nonempty point-image intersection is the image of a
commuting span already in that stage: there are an object $R$ and
arrows $u:R\to P$, $v:R\to Q$ such that
\[
             hu=gv,\qquad (hu)(R)=h(P)\cap g(Q).
\]
Consequently every nonempty finite intersection of arrow images
in one object is itself an existing arrow image.

This statement asserts equality of point sets and the existence
of the displayed factorization span.  It does not assert the
categorical pullback universal property.
\end{proposition}

\begin{proof}
If $h$ factors through $g$, take $R=P$, $u=1_P$, and the given
factor as $v$.  Suppose it does not.  By
Lemma~\ref{C:birth-forbidden-pair}, the nonempty set
\[
                           h^{-1}(g(Q))\subseteq P
\]
contains no forbidden pair.  When $P\ne V$,
Lemma~\ref{C:birth-cover-helly} places that whole set in a single
incoming generating image $e:E\to P$.  Replacing $h$ by $he$
does not change its image intersection with $g$:
\[
                        (he)(E)\cap g(Q)=h(P)\cap g(Q).
\]
The source has moved strictly earlier in its generation history;
for the original object $W$, this step moves to $V$.

Continue until the current arrow factors through $g$.  This must
happen: otherwise the strict descent reaches $V$.  A nonempty
intersection at $V$ says that its point is $g(q)$ for some
underlying point $q\in Q$.  Every point is represented by a
$V$-arrow, so the corresponding map $V\to Q$ gives a
factorization through $g$.  Composing the source-descent maps
produces $u$, and this final factorization produces $v$.

For finite intersections, repeat the two-image assertion.  At
each step the new intersection image is represented by an actual
arrow, so the argument applies again; all earlier factorizations
remain available by composition.
\end{proof}

\begin{proposition}[Proved: pointwise equalizers descend to one arrow]
\label{C:point-equalizer-image}
For every parallel pair $f,g:S\rightrightarrows T$ in a repaired
C.1 stage, its pointwise fixed subset is empty or is exactly the
image of a single arrow $e:R\to S$ already in that stage:
\[
                  \{x\in S:f(x)=g(x)\}=e(R),\qquad fe=ge.
\]
No categorical equalizer universal property is claimed.
\end{proposition}

\begin{proof}
For $f=g$, use $e=1_S$.  Assume $f\ne g$ and that the fixed subset
$F$ is nonempty.  If $F$ contained a forbidden pair in the source
$S$, Lemma~\ref{C:birth-forbidden-pair}, with $h=f$ and comparison
$g$, would factor $f$ through $g$.  That factor would be an
endomorphism of $S$.  All endomorphisms of every constructed
object are identities, so this would force $f=g$, a contradiction.

For $S\ne V$, the Helly lemma therefore places $F$ in one
generating image $e_0:E\to S$.  Its preimage is exactly
$\operatorname{Eq}(fe_0,ge_0)$.  If these two composites are equal,
then $F=e_0(E)$ and the proof is finished.  Otherwise repeat the
argument with the strictly earlier source $E$.

This descent terminates.  At $V$, a nonempty fixed subset means
that the two maps of its singleton agree, and hence are the same
concrete arrow.  Composing the descent maps gives $e:R\to S$.
Injectivity at every descent step gives the asserted exact
equality $F=e(R)$, rather than merely containment.
\end{proof}

\begin{corollary}[Proved: a finite certificate for all future targets]
\label{C:certified-safe-source}
Fix an object $S$ of a repaired C.1 construction.  Suppose every
arrow $e:R\to S$ is injective on binary graph components.  Then
every parallel pair $f,g:S\rightrightarrows T$ at every later
stage has a fixed subgraph connected in each component of $G_S$.
Every later quasi-filter attachment whose source is this $S$
therefore preserves the whole old cycle colimit and the whole
old binary second homology.

The hypothesis is a finite condition which can be checked at the
birth of $S$, provided the retained finite substage contains all
its incoming generators and their ancestors.  No subsequent
stage adds an object with an arrow into $S$, or adds an arrow
between its old objects.
\end{corollary}

\begin{proof}
Proposition~\ref{C:point-equalizer-image} writes any nonempty
fixed subset as $e(R)$ for an actual arrow into $S$.  Its induced
graph is exactly $e_\#G_R$.  The stated hypothesis gives
$\pi_0(F)\to\pi_0(G_S)$ injective, so the exact fixed-path
kernel is zero.  The construction's full-stage and new-maximal-object
rules make the entire incoming category of $S$ permanent from
its birth onward.  It is finite, and the binary component maps
are finite maps, which proves the certificate assertion.
\end{proof}

\begin{corollary}[Proved: the separate component-injectivity obligation]
\label{C:equalizer-component-reduction}
Suppose every arrow of an old repaired C.1 stage induces an
injection on binary graph components.  Then every quasi-filter
attachment for a pair in that stage has zero old cycle-colimit
kernel.  If this component-injectivity assertion holds at every
stage, the first-stage binary second-homology class survives the
whole infinite construction.
\end{corollary}

\begin{proof}
By Proposition~\ref{C:point-equalizer-image}, the fixed subset of
any defining pair is empty or $e(R)$ for an existing arrow
$e:R\to S$.  Arrow images induce their source graphs, by $(c_m)$,
and the equalizer graph is induced on its fixed vertices.
Consequently its graph is precisely $e_\#G_R$.  Component
injectivity of $e$ is exactly the sufficient condition
$\pi_0(F)\to\pi_0(G_S)$ in
Proposition~\ref{C:fixed-path-kernel}.  The fixed-path kernel
therefore vanishes.  Disjoint-sum attachments have zero kernel
as well, and filtered unions preserve the resulting injections
on homology.
\end{proof}

\paragraph{Resolution of the component obligation.}
The descent results above do not by themselves make arrows
injective on binary graph components: a bounded point-transition
relation does not exclude component cycles merely by passage to
components. The simultaneous induction in
Theorem~\ref{C:universal-components} supplies this separate step.
It proves actual-arrow routing for forbidden pairs of components,
represents component-image intersections of size at least two
by older actual spans, and bounds every component path by the
corresponding point-path bound plus one. Consequently
Corollary~\ref{C:equalizer-component-reduction} applies at every
stage, and Corollary~\ref{C:universal-nontotal} completes the
universal-choice assertion in C.1.
